%!TEX root = ../thesis.tex
%*******************************************************************************
%****************************** Fifth Chapter **********************************
%*******************************************************************************

\chapter{Conclusions}\label{chap:conclusions}

% **************************** Define Graphics Path **************************
\ifpdf
    \graphicspath{{Chapter5/Figs/Raster/}{Chapter5/Figs/PDF/}{Chapter5/Figs/}}
\else
    \graphicspath{{Chapter5/Figs/Vector/}{Chapter5/Figs/}}
\fi

\section{Summary of Key Findings}
This thesis investigated the design and feasibility of an agentic retrieval-augmented conversational assistant that supports Greek psychiatric nurses in accessing patient information. Through interviews with practising nurses and the creation of synthetic clinical notes, the study captured the core informational needs of the role of a nurse. The prototype demonstrated satisfactory performance across thirteen evaluation scenarios, accurately recalling medication regimes, appointments, and psychosocial data, while gracefully handling missing information. 
These results indicate that agentic RAG systems can provide timely, context-aware responses in Greek, extending prior work that has largely focused on English-language deployments.


\section{Contributions}
The research makes three primary contributions:
\begin{itemize}
    \item \textbf{Empirical validation of a nurse-facing assistant}: The evaluation confirmed that a conversational interface can reliably serve psychiatric nurses, answering RQ1 and RQ2 by showing faithful retrieval and grounded summarization across representative tasks.
    \item \textbf{Methodological blueprint}: The thesis details a repeatable pipeline that combines qualitative requirements gathering, synthetic data generation, and scenario-based evaluation, offering guidance for future implementations in privacy-constrained clinical environments.
    \item \textbf{Design insights for responsible AI}: By integrating RAG, explicit memory handling, and decline-to-answer behavior, the work provides practical safeguards against hallucinations and inappropriate disclosure, aligning with responsible AI principles discussed in Chapter~\ref{chap:lit-review}.
\end{itemize}

\section{Recommendations for Future Work}
Further research should focus on validating the assistants with real clinical data, enabling assessment of their behavior with authentic patient documentation by real-world practising nurse professionals in real clinical settings. Only then could usability, trustworthiness, and workflow improvements be evaluated, addressing RQ3 more comprehensively. 

Further technical enhancements could include fine-tuning retrievers on actual Greek psychiatric corpora, integrating speech interfaces for intuitive hands-free use, and expanding tools so that patient records can also be updated or external services can be queried in a secured manner. 

Finally, future studies should measure the assistant's impact on nurse workload, burnout, and patient outcomes to quantify its broader benefits.
 
  